\documentclass[12pt,a4paper]{article}

% ─── Pacotes essenciais ───────────────────────────────────────────────────────
\usepackage[utf8]{inputenc}
\usepackage[T1]{fontenc}
\usepackage[brazil]{babel}
\usepackage{helvet}
\renewcommand{\familydefault}{\sfdefault}

% Margens ABNT (NBR 14724)
\usepackage[top=3cm, bottom=2cm, left=3cm, right=2cm]{geometry}

% Espaçamento 1,5 e recuo de parágrafo
\usepackage{setspace}
\onehalfspacing
\usepackage{indentfirst}
\setlength{\parindent}{1.25cm}

% Matemática
\usepackage{amsmath, amssymb, mathtools}
\usepackage{siunitx}
\sisetup{
  output-decimal-marker = {,},
  separate-uncertainty = true,
  locale = DE
}

% Figuras e tabelas
\usepackage{graphicx}
\usepackage{booktabs}
\usepackage{caption}
\usepackage{subcaption}
\captionsetup{font=small, labelfont=bf, justification=centering}

% Seções
\usepackage{titlesec}
\titleformat{\section}{\normalfont\bfseries\MakeUppercase}{\thesection}{1em}{}
\titleformat{\subsection}{\normalfont\bfseries}{\thesubsection}{1em}{}
\titleformat{\subsubsection}{\normalfont\itshape}{\thesubsubsection}{1em}{}

% Referências ABNT
\usepackage[alf]{abntex2cite}

\usepackage[hidelinks]{hyperref}

% ─── Metadados ────────────────────────────────────────────────────────────────
\title{\textbf{Relatório Parcial de Iniciação Científica\\[0.3cm]
Predição de Falhas em Turbinas Eólicas Utilizando Modelos\\
de Aprendizado Profundo e Dados Meteorológicos Exógenos}}
\author{
  Aluno: Thiago Solé Gomes Heleno \\
  Orientador: Prof. Dr. Marcos G. Quiles \\
  Coorientador: Prof. Dr. Mateus Giesbrecht (FEEC/Unicamp) \\[0.3cm]
  Universidade Federal de São Paulo --- UNIFESP \\
  Instituto de Ciência e Tecnologia --- ICT \\
  Bacharelado em Ciência e Tecnologia / Engenharia de Computação
}
\date{\today}

% ─────────────────────────────────────────────────────────────────────────────
\begin{document}

% ─── Folha de Rosto ──────────────────────────────────────────────────────────
\maketitle
\thispagestyle{empty}
\newpage

% ─── Resumo ──────────────────────────────────────────────────────────────────
\section*{Resumo}
% [Resumo do relatório — contexto, objetivo, métodos, resultados parciais]

\vspace{0.5cm}
\noindent\textbf{Palavras-chave:} % [palavras-chave]

\newpage

% ─── Listas ──────────────────────────────────────────────────────────────────
\tableofcontents
\newpage
\listoffigures
\newpage
\listoftables
\newpage

% [LISTA DE ABREVIATURAS E SIGLAS — SCADA, ML, DL, LSTM, CNN, RMSE, etc.]

\setcounter{page}{1}
\pagenumbering{arabic}

% ─── 1 INTRODUÇÃO ────────────────────────────────────────────────────────────
\section{Introdução}
% [Contextualização do tema, motivação, problema de pesquisa]

% ─── 2 JUSTIFICATIVA ─────────────────────────────────────────────────────────
\section{Justificativa}
% [Relevância do projeto: redução de custos O&M, sustentabilidade, lacunas
%   na literatura, contribuição acadêmica]

% ─── 3 OBJETIVO ──────────────────────────────────────────────────────────────
\section{Objetivo}

\subsection{Aprendizado}
% [Conhecer o que são Redes Neurais — objetivo formativo do aluno]

\subsubsection{Fundamentos de Aprendizado de Máquina}
% [Conceitos básicos: regressão, classificação, métricas]

\subsubsection{Redes Neurais Artificiais}
% [Perceptron, MLP, backpropagation]

\subsubsection{Aprendizado Profundo}
% [LSTM, CNN1D, mecanismos de atenção. Acompanhamento de curso externo
%   sobre redes neurais (vídeo-aulas em `aprendizado_curso/`) com
%   exercícios práticos em PyTorch e NumPy/Pandas.]

\subsection{Treinamento}
% [Testar com dados apenas de teste para aprender — experimentação inicial]

\subsubsection{Pipeline de Treinamento}
% [Divisão treino/validação/teste, hiperparâmetros]

\subsubsection{Datasets de Aprendizado}
% [Dados públicos utilizados para fins didáticos]

% ─── 4 FUNDAMENTAÇÃO TEÓRICA ─────────────────────────────────────────────────
\section{Fundamentação Teórica}

\subsection{Manutenção Preditiva em Turbinas Eólicas}

\subsection{Sistemas SCADA}

\subsection{Séries Temporais e Detecção de Anomalias}

\subsection{Arquiteturas de Aprendizado Profundo}

\subsubsection{LSTM}
% [Equações de gates, célula de memória]

% [FIGURA: arquitetura célula LSTM — diagrama gates input/forget/output]

\subsubsection{CNN1D}
% [Convolução 1D, pooling, kernel, stride]

% [FIGURA: arquitetura CNN1D — convolução em janela temporal]

\subsubsection{Mecanismo de Atenção (MultiHeadAttention)}
% [Self-attention, multi-head — usado como camada dentro do CNN-BiLSTM-Attention
%   (NB2). NÃO foi implementada arquitetura Transformer completa. Distinção
%   importante em relação ao plano original.]

% [FIGURA: bloco MultiHeadAttention — Q, K, V, scaled dot-product]

\subsection{Métricas de Avaliação}
% [Definição matemática de Acurácia, Precisão, Recall, F1, AUC-ROC, RMSE, MAE]

\subsubsection{Métricas Padrão de Classificação}
% [Acurácia, Precisão, Revocação, F1-Score, AUC-ROC, AUC-PR, Matriz de Confusão]

\subsubsection{CARE Score (framework adotado)}
% [Composite Anomaly Reliability Evaluation — métrica multi-objetivo
%   composta por F1_2, EF1_2 (event-level F1), WS (Warning System) e
%   acurácia por evento. ATRIBUIÇÃO: framework CARE Score e análise
%   ARCANA foram adotados/replicados do repositório público
%   EnergyFaultDetector (AEFDI):
%   https://github.com/AEFDI/EnergyFaultDetector — não constituem
%   contribuição original deste trabalho.]

% ─── 5 METODOLOGIAS ──────────────────────────────────────────────────────────
\section{Metodologias}

\subsection{Bases de Dados}

\subsubsection{Repositório SCADA Público}
% [Wind Turbine SCADA open data — variáveis, frequência, período]

% [TABELA: descrição das variáveis SCADA — coluna, unidade, faixa]

\subsubsection{Dados Meteorológicos}
% [CPTEC/INPE, OpenWeatherMap — variáveis disponíveis]

\subsubsection{Análise Exploratória dos Dados (EDA)}

% [FIGURA: distribuição das variáveis SCADA — histogramas]

% [FIGURA: matriz de correlação entre variáveis SCADA]

% [FIGURA: série temporal de potência gerada — exemplo de turbina]

% [FIGURA: proporção classe normal vs falha — gráfico de barras mostrando
%   desbalanceamento]

\subsection{Modelos e Notebooks Implementados}

% [Cinco notebooks foram implementados, divididos em dois paradigmas:
%   supervisionado e semi-supervisionado. Tabela resumo abaixo, detalhes
%   nas subseções seguintes.]

% [TABELA: resumo dos 5 notebooks — Notebook | Paradigma | Modelo | Parâmetros]

\subsubsection{NB1 --- CNN-LSTM (Supervisionado, replicação Qi et al. 2024)}
% [Arquiteturas CNN, LSTM e CNN-LSTM (~160K params). Janela 36 timesteps,
%   285 features (XGBoost top 30%). Loss CrossEntropy, Adam lr=1e-3,
%   batch=600, undersampling 1:1 no treino. Threshold ótimo por F1 na
%   validação.]

% [FIGURA: arquitetura CNN-LSTM — Conv1D + MaxPool + LSTM + Dense]

\subsubsection{NB2 --- CNN-BiLSTM-Attention Autoencoder (Semi-supervisionado, PyTorch)}
% [Autoencoder com Conv1D + BiLSTM + camada de atenção MultiHeadAttention
%   + Decoder LSTM (~954K params). Treinado SOMENTE em dados normais.
%   Janela 36 timesteps × 305 features. MAE loss, Adam lr=1e-4, batch=64,
%   AMP. Versão v4 com 5 estratégias pós-hoc: pesos AUROC, suavização
%   temporal, máscara informativa, XGBoost pós-hoc, threshold Beta-F1.
%   Importante: USOU camada de atenção, não Transformer completo.]

% [FIGURA: arquitetura CNN-BiLSTM-Attention AE]

\subsubsection{NB3 --- Dense MLP Autoencoder (Semi-supervisionado, Keras + Optuna)}
% [Autoencoder MLP simétrico (~150K params), bottleneck 35. Sem janela
%   temporal — vetor 305 features por amostra. MSE loss, Adam +
%   ExponentialDecay. Hiperparâmetros otimizados via Optuna (100 trials).
%   Threshold adaptativo (RMSE_pred + gamma). Avaliação CARE Score por
%   evento (58 eventos).]

% [FIGURA: arquitetura Dense MLP AE — encoder/bottleneck/decoder]

\subsubsection{NB4 --- Classifier MLP 3-class com Falhas Induzidas (Supervisionado)}
% [MLP 3-class (Normal/Anomalia Real/Anomalia Induzida). Data augmentation
%   via injeção de 7\% de falhas sintéticas no treino. Hiperparâmetros
%   Optuna: n_layers=2, hidden=192, dropout=0.29, lr=5.47e-3.]

\subsubsection{NB5 --- Autoencoder MLP com Calibração por Falhas Induzidas (Semi-supervisionado)}
% [AE MLP treinado em dados normais, threshold calibrado via falhas
%   sintéticas. Três thresholds comparados: Standard P95, Induced P95,
%   Adaptive (P50_ind+gamma).]

% [FIGURA: diagrama do pipeline geral — entrada → modelo → saída]

\subsection{Data Processing}

\subsubsection{Extração e Conversão}

\subsubsection{Limpeza de Dados e Tratamento de Outliers}

% [FIGURA: boxplot antes/depois da remoção de outliers]

\subsubsection{Imputação de Dados Faltantes}

\subsubsection{Sincronização Temporal}

\subsubsection{Seleção de Variáveis (Feature Selection)}

% [FIGURA: ranking de importância das features]

\subsubsection{Normalização e Escalonamento}

\subsection{Hiperparâmetros e Otimização}
% [Learning rate, batch size, otimizador, épocas, early stopping]

% [TABELA: hiperparâmetros finais por modelo]

\subsection{Estratégia de Validação}
% [Split treino/val/teste, validação cruzada temporal, seeds]

\subsection{Reprodutibilidade}
% [Versões: Python, PyTorch/TensorFlow, sementes aleatórias]

% ─── 6 EXPERIMENTOS ──────────────────────────────────────────────────────────
\section{Experimentos}

\subsection{Configuração Experimental}
% [Hardware: GPUs RTX 4060 / A6000, CINE Cluster ENIAC. Frameworks]

\subsection{Experimentos sem Dados Exógenos}
% [Baseline com apenas dados SCADA]

\subsection{Experimentos com Dados Exógenos}
% [Inclusão de dados meteorológicos — descrever por que tornou-se inviável]

\subsection{Ablation Study}
% [Impacto de cada componente: features, arquitetura, balanceamento]

\subsection{Custo Computacional}
% [Tempo de treinamento, memória GPU consumida, inferência]

% [TABELA: tempo de treinamento e parâmetros por modelo]

% ─── 7 PROBLEMAS COM DADOS ───────────────────────────────────────────────────
\section{Problemas com Dados}

\subsection{Escassez de Dados de Anomalia}
% [Desbalanceamento severo — falhas raras]

% [FIGURA: histograma de classes — normal vs cada tipo de falha]

\subsection{Localização Não Divulgada dos Sítios}
% [Datasets públicos anonimizam coordenadas dos parques eólicos]

\subsection{Impossibilidade de Uso de Dados Meteorológicos Externos}
% [O dataset CARE_To_Compare anonimiza completamente a localização das
%   turbinas (apenas identificadores Wind Farm A/B/C, sem coordenadas
%   geográficas, sem datas absolutas reais). Sem essas informações,
%   é impossível alinhar temporalmente e espacialmente os dados SCADA
%   com bases meteorológicas externas (CPTEC/INPE, OpenWeatherMap).
%   Consequência: a comparação "modelos com vs sem dados exógenos"
%   prevista no plano original não pôde ser executada. O projeto foi
%   readequado para focar na comparação entre paradigmas de detecção
%   (supervisionado vs semi-supervisionado) e na exploração de técnicas
%   alternativas (falhas induzidas como data augmentation).]

\subsection{Qualidade dos Dados SCADA}
% [Ruído, valores faltantes, dessincronização entre sensores]

% [FIGURA: mapa de calor de valores faltantes ao longo do tempo]

% ─── 8 PROBLEMAS NOS MODELOS ─────────────────────────────────────────────────
\section{Problemas nos Modelos}

\subsection{Desbalanceamento de Classes}
% [Estratégias testadas: class weights, oversampling, focal loss]

\subsection{Overfitting}
% [Sintomas observados, regularização aplicada]

% [FIGURA: curvas de loss treino vs validação — evidência de overfitting]

\subsection{Custo Computacional do Treinamento}

\subsection{Generalização para Diferentes Turbinas}

\subsection{Dificuldades de Implementação}
% [Bugs, integração de bibliotecas, limites de hardware]

% ─── 9 COMPARAÇÃO COM OUTRAS PESQUISAS ───────────────────────────────────────
\section{Comparação com Outras Pesquisas}

\subsection{Trabalhos Relacionados na Literatura}
% [Estado da arte — papers sobre detecção falhas turbinas eólicas]

\subsection{Análise Comparativa de Resultados}

% [TABELA: comparação de métricas — nosso modelo vs literatura]

% ─── 10 RESULTADOS OBTIDOS ───────────────────────────────────────────────────
\section{Resultados Obtidos}

\subsection{Resultados Quantitativos}

% [TABELA: métricas consolidadas por notebook — fonte:
%   resultados/metricas_consolidadas.csv. Resumo:
%   NB1 CNN: F1=5.47%, AUC=0.55
%   NB1 LSTM: F1=0% (colapso para classe majoritária)
%   NB1 CNN-LSTM: F1=3.16%, AUC=0.66
%   NB2 P99: F1=5.78%, AUC=0.865
%   NB2 v4 Score Pond.: F1=4.52%, AUC=0.872
%   NB3 Keras AE: AUC=0.872, CARE=0.699, Recall evento=94.1%, Prec evento=88.9%
%   NB4 Induzido: F1=67.4%, Recall=100%
%   NB5 Induced P95: CARE=0.545, Acc=0.98]

% [TABELA: comparação NB3 vs NB5 — pontos opostos da curva ROC,
%   sugere ensemble como linha futura]

\subsection{Resultados Qualitativos}

% [FIGURA: ranking importância XGBoost — JÁ DISPONÍVEL:
%   resultados/02_cnn_bilstm_autoencoder/xgb_feature_importance.png]

% [FIGURA: análise de features (AUROC por feature) — JÁ DISPONÍVEL:
%   resultados/02_cnn_bilstm_autoencoder/feature_analysis.png]

% [FIGURA: curvas ROC e Precision-Recall (NB2) — JÁ DISPONÍVEL:
%   resultados/02_cnn_bilstm_autoencoder/roc_pr_curves.png]

% [FIGURA: comparação induzido vs baseline (NB4) — JÁ DISPONÍVEL:
%   resultados/04_classifier_induced_failure/comparison_plots.png]

% [FIGURA: comparação de thresholds Standard/Induced/Adaptive (NB5) — JÁ DISPONÍVEL:
%   resultados/05_autoencoder_induced_failure/comparison_plots.png]

% [FIGURA A GERAR: matriz de confusão — modelo final por notebook]

% [FIGURA A GERAR: curvas ROC NB1, NB3, NB4, NB5]

% [FIGURA A GERAR: comparação de barras F1-Score entre os 5 NBs]

\subsection{Discussão dos Resultados}

% ─── 11 LIMITAÇÕES DO ESTUDO ─────────────────────────────────────────────────
\section{Limitações do Estudo e Atribuições}

\subsection{Desvios em Relação ao Plano Original}
% [Itens previstos no plano FAPESP que NÃO foram realizados:
%   - Arquitetura Transformer completa (apenas camada de atenção
%     MultiHeadAttention foi utilizada dentro do CNN-BiLSTM-Attention).
%   - Integração com dados meteorológicos exógenos (CPTEC, OpenWeatherMap)
%     — inviável devido à anonimização do dataset CARE_To_Compare.
%   - Comparação direta "modelos com vs sem dados exógenos".]

\subsection{Componentes Replicados de Fontes Externas}
% [Os seguintes componentes foram adotados/replicados de repositórios
%   externos e NÃO constituem contribuição original deste trabalho:
%   - Framework CARE Score (Composite Anomaly Reliability Evaluation)
%   - Análise ARCANA (importância de features)
%   Fonte: repositório público EnergyFaultDetector (AEFDI),
%   https://github.com/AEFDI/EnergyFaultDetector
%   Estes módulos foram utilizados como ferramentas de avaliação
%   padronizadas para permitir comparação com a literatura.]

\subsection{Outras Limitações}
% [Hardware, tempo, escopo do dataset (apenas Wind Farm C),
%   precisão amostral baixa (~2-3\%) inviabiliza uso direto em produção]

% ─── 12 CRONOGRAMA REALIZADO ─────────────────────────────────────────────────
\section{Cronograma Realizado vs Planejado}
% [Comparação A1–A9 do plano original com o executado até o momento]

% [TABELA: cronograma — atividade × bimestre × status (concluído/parcial/atrasado)]

% ─── 13 ATIVIDADES DESENVOLVIDAS ─────────────────────────────────────────────
\section{Atividades Desenvolvidas no Período}
% [Descrição narrativa do que foi feito mês a mês — leitura, código,
%   reuniões, experimentos]

% ─── 14 PRODUÇÃO CIENTÍFICA ──────────────────────────────────────────────────
\section{Produção Científica e Participação em Eventos}
% [Artigos submetidos, apresentações em seminários, congressos.
%   Participação no grupo de pesquisa do Prof. Quiles desde 08/2024]

% ─── 15 MATÉRIAS CURSADAS NA FACULDADE ───────────────────────────────────────
\section{Matérias Cursadas na Faculdade}
% [Disciplinas cursadas no período da IC e relevância para o projeto]

% [TABELA: disciplina | semestre | nota | relação com o projeto]

% ─── 16 LIÇÕES APRENDIDAS ────────────────────────────────────────────────────
\section{Lições Aprendidas}
% [Reflexão pessoal sobre desafios técnicos e metodológicos superados]

% ─── 17 TRABALHOS FUTUROS ────────────────────────────────────────────────────
\section{Trabalhos Futuros}
% [Próximas etapas concretas para o 2º semestre da IC: novos datasets,
%   modelos a explorar, escrita de artigo]

% ─── 18 CONSIDERAÇÕES FINAIS ─────────────────────────────────────────────────
\section{Considerações Finais}
% [Conclusão parcial do trabalho]

% ─── REFERÊNCIAS ─────────────────────────────────────────────────────────────
\bibliographystyle{abntex2-alf}
\bibliography{referencias}

% ─── APÊNDICES ───────────────────────────────────────────────────────────────
\appendix

\section{Apêndice A --- Hiperparâmetros Completos}
% [Lista detalhada de todos hiperparâmetros testados]

\section{Apêndice B --- Trechos de Código Relevantes}
% [Snippets de pré-processamento e arquitetura dos modelos]

\section{Apêndice C --- Resultados Adicionais}
% [Gráficos e tabelas suplementares]

\end{document}
